\documentclass[11pt]{article}
\usepackage[T1]{fontenc}
\usepackage{amsmath,amssymb,amsthm}
\usepackage[margin=1in]{geometry}
\usepackage{booktabs}
\usepackage[hidelinks]{hyperref}
\usepackage{microtype}

\newtheorem{theorem}{Theorem}
\newtheorem{lemma}[theorem]{Lemma}
\newtheorem{proposition}[theorem]{Proposition}
\newcommand{\GB}{\Gamma_{B}}
\newcommand{\qB}{q_{B}}

\title{Nine $B$-Colors Suffice Through Ten Vertices for\\
$3$-Degenerate $K_{3,4}$-Free Graphs of Maximum Degree at Most Four}
\author{Xiaojian Zeng\\
PBC School of Finance, Tsinghua University, 100190\\
\texttt{xiaojianzeng@foxmail.com}}
\date{}
\hypersetup{
  pdftitle={Nine B-Colors Suffice Through Ten Vertices for 3-Degenerate K3,4-Free Graphs of Maximum Degree at Most Four},
  pdfauthor={Xiaojian Zeng},
  pdfkeywords={B-coloring, graph coloring, degeneracy, exact enumeration, SAT, computer-assisted proof}
}

\begin{document}
\maketitle

\begin{abstract}
A $B$-coloring is a proper edge-coloring in which every $4$-cycle is
rainbow.  Hu, Kong, and Wang showed that the maximum $B$-chromatic index
$M_3(4)$ over finite simple $3$-degenerate, $K_{3,4}$-free graphs of maximum
degree at most four is either nine or ten.  We prove the following exact
finite result by exhaustive generation and checked coloring: every graph in
this class on at most ten vertices has $B$-chromatic index at most nine.  The
bound is sharp, since $K_{3,3}$ has $B$-chromatic index nine.  Consequently,
any graph in the class with $B$-chromatic index ten has a connected component
on at least eleven vertices.  The computation enumerates exact isomorphism
classes, checks every returned SAT coloring edge by edge, and separately
repeats each colorability decision with an exact DSATUR backtracker.  This
finite theorem does not determine $M_3(4)$; the original problem remains
unresolved.
\end{abstract}

\noindent\textbf{Keywords.} $B$-coloring; graph coloring; degeneracy;
exact enumeration; SAT; computer-assisted proof.\\
\textbf{2020 Mathematics Subject Classification.} 05C15, 05C35, 05--04.

\section{Problem and contribution}

Let $G$ be a finite simple graph.  A \emph{$B$-coloring} of $G$ is a proper
edge-coloring for which the four edges of every simple $4$-cycle receive
pairwise distinct colors.  Its minimum number of colors is denoted by
$\qB(G)$.  Equivalently, the \emph{$B$-conflict graph} $\GB(G)$ has vertex set
$E(G)$, with two edges adjacent when they share an endpoint or are opposite
edges of a $4$-cycle; then
\[
  \qB(G)=\chi(\GB(G)).
\]
These definitions and the extremal problem considered here are due to Hu,
Kong, and Wang~\cite{HuKongWang2026}.

Call $G$ \emph{legal} if it is $3$-degenerate, has maximum degree at most
four, and contains no copy of $K_{3,4}$ (the copy need not be induced).  Thus
every nonempty subgraph of a legal graph has a vertex of degree at most three.
Hu, Kong, and Wang proved that $\qB(G)\leq 10$ for every legal $G$, while
$K_{3,3}$ gives the lower bound nine.  In their notation,
\[
  9\leq M_3(4)\leq 10,
\]
and they ask which value occurs~\cite[Theorem~2.8 and
Problem~4.1]{HuKongWang2026}.

Our contribution is restricted to graphs of order at most ten.

\begin{theorem}\label{thm:main}
Every finite simple legal graph $G$ on at most ten vertices satisfies
$\qB(G)\leq 9$.  This bound is sharp.  Consequently, if a legal graph $H$
satisfies $\qB(H)=10$, then some connected component of $H$ has at least
eleven vertices.
\end{theorem}

This is an order-bounded theorem, not a resolution of the extremal problem.
Exact finite certification has also recently been used for restricted graph
problems in which a structural reduction is followed by canonical cases and
independently checked witnesses; for example, Zeng~\cite{Zeng2026Tuza}
reduces a split-graph statement to $44{,}702$ canonical records carrying
explicit packing and covering certificates.  Here the objects are instead
unlabeled legal graphs, and the decisive witness for each class is a verified
nine-coloring of its $B$-conflict graph.  This comparison is methodological;
neither result implies the other.

\section{Reduction to connected isomorphism classes}

\begin{lemma}\label{lem:components}
If $G$ has connected components $G_1,\dots,G_t$, then
\[
  \GB(G)=\GB(G_1)\mathbin{\dot\cup}\cdots\mathbin{\dot\cup}\GB(G_t)
  \quad\text{and}\quad
  \qB(G)=\max_i\qB(G_i).
\]
\end{lemma}

\begin{proof}
Two edges in distinct components neither meet nor lie on a common $4$-cycle.
Thus there are no conflict-graph edges between different components, and the
chromatic number of the disjoint union is the maximum of the component
chromatic numbers.
\end{proof}

It therefore suffices to enumerate connected legal graphs.  Orders at most
seven are taken from the Graph Atlas of Read and Wilson~\cite{ReadWilson1998}
through NetworkX's \texttt{graph\_atlas\_g} interface
\cite{HagbergSchultSwart2008}.  That interface contains one representative
of every simple isomorphism class on at most seven vertices.

For the subsequent orders we use the following elementary induction.

\begin{lemma}\label{lem:generation}
Suppose one representative of every connected legal graph on $n-1$ vertices
is available, where $n\geq 2$.  Extending every representative by a new
nonisolated vertex in every neighborhood of size at most four, discarding
illegal graphs, and deduplicating by exact graph isomorphism produces one
representative of every connected legal graph on $n$ vertices.
\end{lemma}

\begin{proof}
Every connected graph on at least two vertices has a vertex whose deletion
leaves the graph connected: take a leaf of a spanning tree.  Maximum degree
at most four, $3$-degeneracy, and $K_{3,4}$-freeness are hereditary under
vertex deletion.  Deleting such a vertex from a connected legal $n$-vertex
graph therefore gives a connected legal graph represented at order $n-1$.
Restoring the vertex is precisely one of the enumerated nonempty
neighborhoods, and each chosen neighbor had degree below four before the
restoration.  Conversely, the legality test discards every extension outside
the class.  Exact isomorphism deduplication proves the assertion.
\end{proof}

The implementation first buckets extensions by degree sequence, triangle
count, and a Weisfeiler--Lehman hash.  These quantities are used only as
isomorphism invariants to reduce comparisons.  Every decision that two graphs
are duplicates is made by an exact NetworkX graph-isomorphism test; the hash
is never treated as a certificate of isomorphism.

The legality tests directly implement the defining predicates.  Maximum
degree is read from the degree sequence.  A graph is $3$-degenerate exactly
when its maximum core number is at most three.  Finally, a simple graph
contains a (not necessarily induced) $K_{3,4}$ exactly when some triple of
vertices has at least four common neighbors.  At order ten an optimized
extension test is used.  Since the base is legal, any newly created
$K_{3,4}$ uses the new vertex.  If the new vertex lies on the four-vertex
side, three vertices in its neighborhood have three old common neighbors.  If
it lies on the three-vertex side, its degree is four and two old vertices are
adjacent to all four of its neighbors.  These are exactly the two cases
tested.  Before the full order-ten run, this optimized predicate was compared
with the complete common-neighbor predicate on every extension of the first
$300$ order-nine representatives, with no disagreement.

\section{Exact coloring certification}

For each representative $G$, the program orders $E(G)$ and constructs
$\GB(G)$.  For distinct edges $uv$ and $xy$, it inserts a conflict when the
edges meet or, if they are disjoint, when either
\[
  ux,vy\in E(G) \qquad\text{or}\qquad uy,vx\in E(G).
\]
The latter condition is equivalent to the two edges being opposite on a
simple $4$-cycle.

Nine-colorability is encoded in CNF.  For every vertex $e$ of $\GB(G)$ and
color $c\in\{1,\dots,9\}$ there is a Boolean variable $x_{e,c}$.  Clauses
require exactly one color at each vertex, and for every conflict $ef$ and
every color $c$ the clause $\neg x_{e,c}\lor\neg x_{f,c}$ forbids equality.
The instances are solved through PySAT~\cite{IgnatievMorgadoMarquesSilva2018}
using Glucose4, from the Glucose family of CDCL solvers
\cite{AudemardSimon2009}.  Every satisfying assignment is decoded, checked
to give exactly one color per conflict-graph vertex, and checked against every
conflict edge.

As a second colorability decision, a separate exact backtracker uses the
DSATUR rule~\cite{Brelaz1979}: it branches on an uncolored vertex with maximum
color saturation, breaking ties by degree, and exhausts all currently
available colors.  This route shares the enumerated graphs and conflict-graph
construction with the SAT route, but not its encoding or solver.  The two
decisions agreed for every enumerated class.

\begin{table}[ht]
\centering
\caption{Exact exhaustive screens.  The first row combines all orders at
most seven; raw extensions and isomorphism comparisons apply only to the
inductively generated top order.}
\label{tab:counts}
\begin{tabular}{@{}rrrrr@{}}
\toprule
Order & Connected legal classes & Raw legal extensions & Exact iso comparisons
  & Non-$9$-colorable \\
\midrule
$\leq 7$ & $456$    & ---       & ---       & $0$ \\
$8$       & $1{,}918$  & $15{,}791$  & $13{,}929$  & $0$ \\
$9$       & $12{,}177$ & $116{,}290$ & $104{,}581$ & $0$ \\
$10$      & $89{,}299$ & $933{,}827$ & $849{,}829$ & $0$ \\
\bottomrule
\end{tabular}
\end{table}

All counts in Table~\ref{tab:counts} are integer counts of exact objects, and
all colorability answers are exact.  No floating-point calculation is used
to decide legality, isomorphism, conflict, or colorability.

\begin{proof}[Proof of Theorem~\ref{thm:main}]
By Lemma~\ref{lem:generation}, the Atlas screen and the three inductive
generation steps cover every connected legal isomorphism class on at most ten
vertices.  The checked SAT coloring gives a proper nine-coloring of
$\GB(G)$ for every representative; isomorphism preserves this property.
Lemma~\ref{lem:components} extends the conclusion to arbitrary legal graphs
on at most ten vertices; the empty graph is immediate.

For sharpness, $K_{3,3}$ is legal and has nine edges.  Any two of its edges
either meet or are opposite on a $4$-cycle, so $\GB(K_{3,3})=K_9$ and
$\qB(K_{3,3})=9$.  Finally, if a legal graph has $B$-chromatic index ten,
Lemma~\ref{lem:components} gives a component of index ten.  The finite result
forces that component to have at least eleven vertices.
\end{proof}

\section{Reproducibility and limitations}

The accompanying source archive contains four Python programs and their
recorded outputs: \texttt{atlas\_screen.py},
\texttt{n8\_sat\_screen.py}, \texttt{n9\_sat\_screen.py}, and
\texttt{n10\_sat\_screen.py}.  The top-order program reconstructs all lower
orders before generating order ten.  The recorded full run used Python
$3.10.12$, NetworkX $3.4.2$, PySAT $1.8.dev24$, and the PySAT solver name
\texttt{g4}.  A clean full rerun reproduced the order-ten counts and all
recorded deterministic digests.  The programs also record graph6 digests of
the representative multisets and digests of the checked coloring records,
allowing byte-level comparison of future runs.

The theorem is finite.  In particular, the absence of a ten-color witness
through order ten cannot be extrapolated to graphs of arbitrary order.  No
finite-order reduction beyond Lemma~\ref{lem:generation} is claimed, and no
universal nine-coloring argument is supplied.  Thus the possibilities
$M_3(4)=9$ and $M_3(4)=10$ both remain compatible with this result.

\paragraph{Computational and AI assistance.}
OpenAI Codex assisted with route exploration, implementation and checking of
the enumeration, literature triage, and preparation of the manuscript.  SAT
solving was used only to produce colorings that were checked directly; the
separate DSATUR backtracker repeated every decision.  The author reviewed the
mathematical statement, proof, code, computational scope, citations, and
manuscript, and takes responsibility for them.  No AI system was treated as
an author or as mathematical authority, and no claim of formal proof-assistant
verification is made.

\bibliographystyle{plain}
\bibliography{references}

\end{document}
